\section{Mechanism: Causal Attribution via Embedding-Source Interpolation}
\label{sec:mechanism}

% Full grounding: DELTANET_RD_EXACTNESS_DESIGN.md \S3, \S4, \S7 (tiers);
% EXPERIMENT_LOG.md "Wave 1 ATTRIBUTION VERDICT" 2026-07-04.

\subsection{Design}
\begin{itemize}
  \item Hold grammar/model/loss FIXED, vary ONLY the embedding source
    feeding $k_{\mathrm{eff}}$: learned (baseline, arm iii), frozen
    orthonormal (i), frozen real GPT-2 span embeddings (ii), Gram-matched
    frozen non-orthonormal control (iv).
  \item Claim tier: ``causal, premise-conditional.'' Bundled-vs-geometry-
    only distinction is load-bearing: frozen-vs-learned (i/ii/iv vs.\ iii)
    bundles geometry AND trainable-param count; (i)-vs-(iv) is the
    GEOMETRY-ONLY contrast (conditional on iv's re-verification gate).
\end{itemize}

\subsection{Result: every geometry converges to the same attractor (Fig.~4)}
\begin{itemize}
  \item Arms i/ii/iv ALL converge to the learned baseline's $k_{\mathrm{eff}}$
    geometry: gram dev $2.71$--$2.92$ at $K{=}32$ vs.\ learned
    $2.71$--$2.77$.
  \item Raw embedding geometry is causally IRRELEVANT --- the trained
    $W_k$/conv path maps every input geometry onto the SAME
    non-orthonormal write-geometry basin.
\end{itemize}

\subsection{Complementary evidence: measured-$\beta$ reconstruction}
\begin{itemize}
  \item Closed-form crosstalk reconstruction, measured $\beta$: residual
    $0.0038$ ($K{=}16$) / $0.0041$ ($K{=}32$) vs.\ $0.148/0.377$ under the
    wrong-regime $\beta{=}1$ idealization.
  \item ``Geometry $+\beta$'' is essentially the complete state-formation
    account --- \textbf{at the weaker, non-causal explanatory tier}: Test
    A (idealization adequacy) earns NO explanatory claim; Test B
    (predicted-vs-measured frontier match, e.g.\ $K{=}32$: 0.726 predicted
    vs.\ 0.789 measured) earns ``structural, deterministic reproduction,''
    explicitly NOT a causal claim. State this tier distinction crisply ---
    exactly the kind of blur a reviewer will probe.
\end{itemize}

\subsection{The existence proof: i-strong (Fig.~5)}
\begin{itemize}
  \item Non-trainable, hand-built, per-identity $k_{\mathrm{eff}}$ pin,
    bypassing $W_k$/conv entirely.
  \item Gram dev $0.000$; $\mathrm{rec@0.9} = 1.00/1.00/1.00$ at
    $h{=}1/2/3$, $K{=}32$ --- matches the SYNTHETIC harness exactly.
  \item Claim tier: ``diagnostic tie-breaker,'' NEVER a standalone causal
    headline (two-variable by design: embedding source AND projection
    bypass, bundled).
  \item The paper's turning point: from ``here is a gap'' to ``here is
    what would close it.'' The architecture and the delta rule's own
    algebra CAN represent and use the perfect solution; ordinary training
    does not find it.
\end{itemize}

\subsection{The 2$\times$2 existence proof: orthogonality and stability are
each necessary, jointly sufficient (Fig.~5)}
\begin{itemize}
  \item A 2$\times$2 grid cutting the same archive along two independent
    geometric axes (orthogonal $\times$ stable), $K{=}32$: stable+orthogonal
    (i-strong, pool-restricted to $32{\leq}d_{\mathrm{state}}$) $h{=}4$
    $1.0000$, $h{=}7$ $1.0000$; stable+not-orthogonal (frozen-embedding arms
    i/iv) $h{=}4$ $0.0075$--$0.0114$, $h{=}7$ ${\approx}0.0000$;
    not-stable+orthogonal (bare geo3 alone) $h{=}4$ $0.44$ mean, $h{=}7$
    $0.018$ mean; not-stable+not-orthogonal (fully-learned baseline)
    $h{=}4$ $0.009$, $h{=}7$ ${\approx}0.0000$.
  \item Neither axis alone clears the $\sim$0.01 floor; orthogonality alone
    recovers most of the gap ($44$--$56\times$) but plateaus at $0.44$;
    adding stability is not a partial improvement over that $0.44$, it is a
    discontinuous jump to $1.00$ --- the INTERACTION, not either main
    effect, is the figure's point.
\end{itemize}

\subsection{Closing this section: whether SGD can reach the ``jointly
sufficient'' corner itself --- the key-anchoring program's complete arc
(\S\ref{sec:the-fix}, \S\ref{sec:results}, Fig.~10, Fig.~11)}
\begin{itemize}
  \item The 2$\times$2's ``jointly sufficient'' corner was, until this
    program ran, demonstrated only by the hand-built i-strong pin. The
    natural question --- can a LEARNED mechanism reach that corner via
    SGD? --- was tested to completion across four waves, and the answer is
    now final: \textbf{no, not via the hypothesized channel, but the
    program found something arguably more interesting.}
  \item \textbf{The behavioral gain is real and independently replicated.}
    A learned, per-entity cross-episode key anchor lifts $K{=}32$ $h{=}4$
    held-out recovery from a freshly-measured baseline $0.4105$ to
    $0.61$--$0.71$ (9 seed-runs across 3 distinct seed sets, 2 architecture
    variants).
  \item \textbf{The hypothesized entity-alignment mechanism is REJECTED
    (Outcome C).} A purpose-built per-entity engagement instrument
    (\texttt{engaged\_frac\_v3}, median $r_e$) shows fewer than
    $14$--$41\%$ of entities individually engage at every measured leg,
    confirmed across two independent mechanism-tier waves at increasing
    evidentiary rigor (unit-anchor-norm measured directly, a formal
    dip-test ruling out masked bimodality, an effect-size floor ruling out
    power-manufactured engagement), immune to a specific rescue attempt
    (a $\lambda$-scaled threshold re-derivation) that was itself
    independently attacked and rejected.
  \item \textbf{The actual explanation was found and CONFIRMED BY
    ABLATION.} A frozen, never-trained anchor table (random init or
    frame-potential init, both tested) at the same blend weight matches or
    exceeds the learned table's own gain ($0.67$--$0.76$ vs.\
    $0.61$--$0.71$), with a negative-median $r_e$ reading in the
    random-init arm confirming the instrument correctly reports ``nothing
    to align to.'' Constancy in the blend arithmetic --- not learned
    entity-specific content --- is the active ingredient; SGD's only
    demonstrated contribution is tuning the blend weight.
  \item \textbf{This fix's own regime is now bounded, not open-ended ---
    and the boundary is now LOCATED, then shown to DISSOLVE with scale
    (Fig.~cliff).} Extending the identical design to $K{=}48$ ($K/d{=}0.75$)
    shows both named explanatory channels measurably active yet the fix
    still misses its own bar 0/3 seeds --- the capacity curve completes as
    a cliff ($\sim$1.00 $\to$ $\sim$0.65 $\to$ $\sim$0.02 at
    $K/d{=}0.25/0.5/0.75$), not a smooth decline, while the theoretical
    ceiling itself declines only gently over the same range. A dedicated
    follow-on localizes this transition's midpoint to $x_0{=}0.5455$ (95\%
    CI $[0.5385,0.5513]$) at $d_{\mathrm{state}}{=}64$, then shows it is
    NOT a universal property of $K/d$: the identical ratio window at
    $d_{\mathrm{state}}{=}128$ is flat at $h4{=}1.0$ across all four
    matched loads, 12/12 cells --- the transition exits the window
    entirely rather than merely shifting. Directly injecting controlled,
    frozen anchor-table coherence (the leading confound between the two
    scales, forced by the Welch bound only when entity count exceeds
    $d_{\mathrm{state}}$) up to and beyond $d{=}64$'s own realized training
    band, under a concentrated rank-4 structure, produces no cliff at any
    dose (Fig.~dose) --- exonerating scalar coherence as a sufficient
    account and leaving overlap structure and raw state capacity as the
    two live, untested candidates (\S\ref{sec:discussion}). \textbf{The
    raw-state-capacity candidate has since been tested directly and load
    tolerance turns out to grow SUPER-LINEARLY in $d_{\mathrm{state}}$,
    not at a fixed $K/d$ ratio:} re-fitting the transition at
    $d_{\mathrm{state}}{=}80$ places it at $x_0{=}0.6779$ (95\% CI
    $[0.6683,0.6867]$), cleanly excluding the ratio-invariant band around
    $0.5455$; at $d_{\mathrm{state}}{=}96$ the FULL unlocked grid
    ($K{\in}\{69,72,78,84,90\}$, $n{=}3$/$K$ --- 11 of 12 new cells were
    initially flagged inadmissible by a since-diagnosed tolerance
    miscalibration in the convergence gate, not a real capacity wall; all
    4,608 flagged episodes resolve by $n_{\mathrm{iter}}{\leq}28$, and only
    the admission flag was recalibrated) shows NO cliff through
    $K/d{=}0.9375$ but a non-monotonic curve
    ($0.9592/0.9216/0.9326/0.9581/1.0000$) that the sigmoid fit cannot
    localize (100\% bootstrap-degenerate). The exact functional form and
    the precise $d{=}96$ transition point remain open --- neither the
    absolute-slack band $[0.718,0.739]$ nor the power-law band
    $[0.768,0.837]$ is confirmed or excluded; seed escalation at the
    noisiest group ($K{=}72$) is the registered next step
    (\S\ref{sec:discussion}).
  \item \textbf{State this precisely, not as a substitution:} the 2$\times$2's
    ``jointly sufficient'' corner still rests on the hand-built i-strong
    pin for a \emph{causal} demonstration at $K{=}32$; the \emph{practical},
    SGD-adjacent route tested to reach that corner does not use the
    hypothesized mechanism at all, works via a simpler (constancy-only)
    route instead, and that alternate route's own working regime is now
    known (holds at $K/d{\leq}0.5$, fails at $K/d{=}0.75$). Both facts
    belong in the paper, clearly separated, neither substituting for the
    other.
\end{itemize}
